% CAPÍTULO 1 - INTRODUÇÃO
%
% Edite o conteúdo abaixo com seu texto. Remova os marcadores [colchetes].
% Use \section{} para seções e \subsection{} para subseções.
% Use \cite{chave} para citações e \ref{label} para referências cruzadas.

\chapter{INTRODUÇÃO}
\label{cap:introducao}

[Texto de contextualização e apresentação do tema da pesquisa.
Exemplo de citação no formato autor-ano: \cite{exemplo2024}.
Exemplo de citação inline: segundo \citeonline{exemplo2024}, o uso de
modelos computacionais tem se mostrado promissor.]

[Texto de justificativa e relevância do trabalho.]

[Texto de apresentação do problema e lacuna de pesquisa.]

[Texto de apresentação da abordagem e contribuição pretendida.]

% -----------------------------------------------------------------------------
\section{Objetivos}
\label{sec:objetivos}

\subsection{Objetivo geral}
\label{sec:objetivo-geral}

[Descreva o objetivo geral do trabalho.]

\subsection{Objetivos específicos}
\label{sec:objetivos-especificos}

\begin{itemize}
    \item {}[Primeiro objetivo específico.]
    \item {}[Segundo objetivo específico.]
    \item {}[Terceiro objetivo específico.]
    \item {}[Quarto objetivo específico.]
\end{itemize}

% -----------------------------------------------------------------------------
\section{Estrutura do trabalho}
\label{sec:estrutura}

[Descreva brevemente como o trabalho está organizado, mencionando
os capítulos e seus respectivos conteúdos. Use \texttt{\textbackslash ref}
para referenciar capítulos, por exemplo: Capítulo~\ref{cap:fundamentacao}.]
